\documentclass[10pt,a4paper]{amsart}
\usepackage[margin=19mm]{geometry}
\usepackage{amsmath,amssymb,mathtools}
\usepackage[scr=rsfs,cal=euler]{mathalfa}
\usepackage{xfrac}
\usepackage[unicode,hidelinks,bookmarks=false]{hyperref}
\newcommand{\ie}{i.e.\ }
\newcommand{\cf}{cf.\ }
\newcommand{\resp}{respectively\ }
\newcommand{\Subsection}{\subsection}
\providecommand{\nobreakdash}{\nobreak\hbox{-}}
\setlength{\emergencystretch}{2em}
\multlinegap0pt
\usepackage{amssymb}
\usepackage{mathtools}
\usepackage{enumerate}
\newtheorem{lemma}{Lemma}
\title{Complete subvarieties in the projectivized strata of meromorphic differentials}
\author{Dawei Chen, Guillaume Tahar}
\date{Source: arXiv:2511.07102v3 (CC BY 4.0)}
\begin{document}
\maketitle

\noindent\emph{Source and licence.} Complete subvarieties in the projectivized strata of meromorphic differentials. Version arXiv:2511.07102v3 (CC BY 4.0), distributed by arXiv under the Creative Commons Attribution 4.0 International licence, http://creativecommons.org/licenses/by/4.0/, as recorded in the arXiv licence metadata for this exact version and shown on the abstract page https://arxiv.org/abs/2511.07102v3. Full licence notice: \texttt{licenses/arxiv-2511.07102-CC-BY-4.0.txt}.\par\smallskip

\noindent\textit{Editorial scope.} Counted unit: the article's lemma lem:Persistence (source0.tex lines 233--235). The article's complete original proof of it is delivered with it (source0.tex lines 237--261, one \texttt{proof} environment of 25 lines). No mathematics is written, repaired or re-derived: the statement and the proof are the article's own lines, in the article's own notation, and no retained line is substituted or re-flowed.  Retained uncounted as the source-local context the proof needs: lines 263--266.  One declared adaptation: exactly 1 ancillary strings inside the retained spans are omitted (image-inclusion commands and, where the source repeats a label, the repeated label definitions) and no other line differs from the source; the figure environments, with their placement, captions and labels, are kept, and the package records the omitted strings in fidelity receipts bound to the affected bodies.  No retained line contains a citation, so the wrapper carries no bibliography and no bibliography entry.  The retained lines contain the article's own diagram code, which the wrapper typesets with the standard tikz packages; no image file is delivered and the package declares no asset dependency.  Editorial formatting only, in addition: an AMS \texttt{amsart} preamble that restates the article's own declarations and macros used by the retained lines, namely the environments \texttt{lemma} on separate counters, together with the standard packages those lines need.  The retained environments print in this document as Lemma 1 in order of appearance, whereas the article prints the same items with the numbering of its own full document.  The complete native source of the article as distributed in the arXiv e-print for this exact version is bundled unchanged as \texttt{sources/188780/source0.tex}.
\begin{figure}

\caption{Deformation of a chain of peripheral saddle connections}\label{Figure:Deformation}
\end{figure}

\begin{lemma}\label{lem:Persistence}
Given a translation surface $(X,\omega)$ in a stratum $\Omega\mathcal{M}_{g}(\mu)$ of meromorphic differentials, every coherent homology class of $(X,\omega)$ persists as a coherent homology class in a neighborhood of $(X,\omega)$ in $\Omega\mathcal{M}_{g}(\mu)$.
\end{lemma}

\begin{proof}
Let $[\delta]=\sum\limits_{\gamma \in E} [\gamma]$ be a coherent homology class of $(X, \omega)$. In particular, there exists $\theta \in \mathbb{R}$ such that for any class $[\gamma] \in E$, $\Re\left(e^{-i\theta} P_{\gamma} \right)>0$.
\par
In the boundary of any polar domain, along an arbitrarily small deformation, a peripheral saddle connection $\gamma$ representing $[\gamma]$ whose incident corners have angle strictly larger than $\pi$ persists as a peripheral saddle connection. Moreover, $\Re\left(e^{-i\theta} P_{\gamma} \right)$ stays positive.
\par
In contrast, given some polar domain, consider a pair of consecutive boundary oriented saddle connections $\gamma_{1},\gamma_{2}$ such that:
\begin{itemize}
    \item $\gamma_{1}$ and $\gamma_{2}$ meet at a boundary corner of angle $\pi$;
    \item the starting point of $\gamma_{1}$ and the ending point of $\gamma_{2}$ are corners of angle strictly larger than $\pi$.
\end{itemize}
We have that $[\gamma_{1}]$ and $[\gamma_{2}]$ belong to $D_\omega$. Assuming that they belong to $E$, $\Re\left(e^{-i\theta} P_{\gamma_{1}} \right)$ and $\Re\left(e^{-i\theta} P_{\gamma_{2}} \right)$ stay positive for any small enough deformation. If along an arbitrarily small deformation, the angle at the corner between $\gamma_{1}$ and $\gamma_{2}$ becomes strictly smaller than $\pi$, then $\gamma_{1}$ and $\gamma_{2}$ cease to be peripheral saddle connections for this polar domain. However, $\gamma_{1}$ and $\gamma_{2}$ form a triangle with a third saddle connection $\gamma'$ satisfying $[\gamma']=[\gamma_{1}]+[\gamma_{2}]$ (the assumption that boundary edges of the same polar domain have counterclockwise orientation is crucial here). This new saddle connection $\gamma'$ is a peripheral saddle connection for the same polar domain and $\Re\left(e^{-i\theta} P_{\gamma'} \right)= \Re\left(e^{-i\theta} P_{\gamma_{1}} \right) +\Re\left(e^{-i\theta} P_{\gamma_{2}} \right)$ stays positive provided that the deformation is small enough. In this case, $[\gamma_{1}]$ and $[\gamma_{2}]$ disappear from $E$ but they are replaced by $[\gamma']$ so that the homology class $[\delta]=\sum\limits_{\gamma \in E} [\gamma]$  does not change.
\par
More generally, consider a chain of consecutive boundary oriented saddle connections $\gamma_{1},\dots,\gamma_{k}$ meeting at corners of angle $\pi$ and such that $[\gamma_{1}],\dots,[\gamma_{k}]$ belong to $E$. Under an arbitrarily small deformation, some of these saddle connections may cease to be peripheral. However, they are replaced by newly created peripheral saddle connections in such a way that the homology class $[\delta]=\sum\limits_{\gamma \in E} [\gamma]$ remains unchanged, see Figure~\ref{Figure:Deformation}.
\par
In the specific case where the chain of consecutive saddle connections is cyclic, it forms the boundary of a polar domain consisting of a semi-infinite translation cylinder foliated by a family of parallel closed geodesics. Denoting by $\alpha_{1}(\varepsilon),\dots,\alpha_{k}(\varepsilon)$ the sequence of corner angles along a deformation parametrized by $\varepsilon$, the total angle defect
$
\sum\limits_{i=1}^{k}\bigl(\pi-\alpha_i(\varepsilon)\bigr)
$
can be interpreted as the winding number of a loop homotopic to a waist curve of the cylinder, and is therefore identically equal to $0$ throughout the deformation. Consequently, the corner angles cannot all become strictly smaller than $\pi$, so the previous arguments still apply, and the boundary of the polar domain continues to be formed by peripheral saddle connections whose homology classes belong to $E$.
 In all cases,
$
[\delta]=\sum_{\gamma\in E}[\gamma]
$
persists as a coherent homology class under sufficiently small deformations.
\end{proof}

\end{document}
